\documentclass[11pt,a4paper]{article}
\usepackage[utf8]{inputenc}
\usepackage[T1]{fontenc}
\usepackage[french]{babel}
\usepackage{geometry}
\usepackage{xcolor}
\usepackage{hyperref}
\usepackage{titlesec}

\geometry{margin=1in}
\definecolor{titleblue}{RGB}{0,51,102}

\hypersetup{
    colorlinks=true,
    linkcolor=titleblue,
    urlcolor=titleblue,
    pdftitle={Lettre de Motivation - Ismail AISSAOUI},
}

\begin{document}

\noindent \textbf{Ismail AISSAOUI} \\
Ingénieur d'État en Intelligence Artificielle (ESI-SBA) \\
Email: ismail.aissaoui.pro@gmail.com \\
Téléphone: +213 660 70 77 96 \\
LinkedIn: https://ismail-aissaoui.vercel.app

\bigskip

\begin{flushright}
    Au Comité de Sélection du Doctorat \\
    \textbf{Programme EDT - Université Côte d'Azur} \\
    À l'attention de : Dr. Gérald Rocher \\
    \medskip
    1er mai 2026
\end{flushright}

\bigskip

\noindent \textbf{Objet : Candidature au doctorat : Raisonnement abductif distribué pour des Jumeaux Numériques auto-explicatifs (Rang 3 - PC1)}

\bigskip

Monsieur le Docteur,

Ingénieur d'État en Intelligence Artificielle diplômé de l'ESI-SBA, je vous adresse avec enthousiasme ma candidature pour le poste de doctorat intitulé « Raisonnement abductif distribué pour des Jumeaux Numériques auto-explicatifs » au sein du programme national Engineering Digital Twin (EDT). Cette candidature est motivée par la pertinence de mon expertise en IA distribuée pour vos recherches sur l'explicabilité hiérarchique.

Actuellement en fin de cycle ingénieur chez Sonatrach, je travaille sur le déploiement de modèles de substitution haute performance (**Mamba-SSM** et **xLSTM**) sur des dispositifs industriels en périphérie (Edge). Cette expérience m'a confronté à la nécessité critique de fournir non seulement des prédictions précises, mais aussi des explications compréhensibles pour les opérateurs de terrain. Votre proposition d'utiliser le raisonnement abductif pour automatiser l'explication des comportements de modèles « boîte noire » dans des architectures distribuées Edge-Fog-Cloud correspond parfaitement à mes aspirations scientifiques.

Je suis particulièrement attiré par l'approche hybride combinant IA symbolique et connexionniste. Ma maîtrise de PyTorch et mon expérience dans l'optimisation de modèles pour le matériel contraint me permettront de contribuer activement à l'implémentation de vos algorithmes de raisonnement au sein de la **plateforme Artemis**.

Travailler au sein de l'Université Côte d'Azur et du consortium EDT représente pour moi une opportunité exceptionnelle de participer à la création de la prochaine génération de jumeaux numériques transparents et fiables.

Je vous remercie de l'attention que vous porterez à ma candidature et reste à votre entière disposition pour un entretien.

Dans l'attente de votre réponse, je vous prie d'agréer, Monsieur le Docteur, l'expression de mes salutations distinguées.

\bigskip

\noindent Ismail AISSAOUI

\end{document}
